\documentclass[12pt]{article}
\usepackage[utf8]{inputenc}
\usepackage[T1]{fontenc}
\usepackage{hyperref}
\title{Sri Lankan Malay Creole Language Resources for K-12 Curriculum Development}
\author{Compiled by Claude Sonnet 5.5 (claude-sonnet-5-5) for Andr\'as Kornai}
\date{}
\begin{document}
\maketitle

\section{Introduction}
Sri Lankan Malay (ISO 639-3 sci; Glottocode sril1245) is a Malay-based creole spoken by the Sri Lankan Malay community, with strong Sinhala and Tamil influence. It is endangered. Other names include Sri Lankan Creole Malay and Melayu Bahasa. Searches found scholarly descriptions and some video documentation, but no school materials.

\section{Primary Educational Resources}
No primers, textbooks or school programmes for Sri Lankan Malay were found. School instruction in Sri Lanka is mainly in Sinhala and Tamil. Researchers report that some community members favor teaching the language, while others favor teaching standard Malaysian or Indonesian Malay instead.

\section{Linguistic and Reference Materials}
\begin{itemize}
\item Smith and Paauw, ``Sri Lanka Malay'' (Creole Language Library vol.\ 29), on structure and Tamil influence. \url{https://www.degruyterbrill.com/document/doi/10.1075/cll.29.09smi/html}
\item DoBeS documentation project page, Sri Lanka Malay (Max Planck Institute for Psycholinguistics). \url{https://dobes.mpi.nl/?p=1300}
\item Endangered Languages Project resource list, which includes ``A Grammar of Upcountry Sri Lanka Malay'', ``Sri Lankan Malay: A Unique Creole'' and ``The genesis of Sri Lanka Malay as a multi-layered process''. \url{https://endangeredlanguages.com/node/38970/resources}
\item Wikipedia overview. \url{https://wikipedia.com/wiki/ISO_639:sci}
\end{itemize}
No dictionary was found in my searches.

\section{Religious and Community Resources}
No scripture translation or community organization specific to the language was found in search results.

\section{Digital and Audio Resources}
\begin{itemize}
\item Endangered Languages Project video of two speakers (Zeena and Meluri) discussing the future of the language. \url{https://www.youtube.com/watch?v=3KGABYk7mg0}
\item A second ELP video (Putresia and Noor, August 2025), listed on the same resource page. \url{https://endangeredlanguages.com/node/38970/resources}
\end{itemize}

\section{Curriculum Development Context}
The language is not known to be taught in schools. Transmission to young people is weakening, partly because English is chosen at home. A curriculum would first need a community-agreed writing system, a lexicon and recorded texts. The DoBeS corpus and published grammars are the best starting points for such work.
\end{document}